\documentclass[11pt]{article}
\usepackage[margin=1in]{geometry}
\usepackage[T1]{fontenc}
\usepackage{lmodern}
\usepackage{amsmath,amssymb,bm,booktabs,tabularx,array}
\usepackage{graphicx,float,placeins,microtype,fancyhdr}
\usepackage[hidelinks]{hyperref}
\usepackage{xurl}
\hypersetup{pdftitle={HPS GPR v6.2: MC injection and signal recovery},pdfauthor={Emrys Peets}}
\pagestyle{fancy}\fancyhf{}\setlength{\headheight}{14pt}
\fancyhead[L]{\small HPS Gaussian-Process Resonance Search}
\fancyhead[R]{\small v6.2 / Signal recovery}\fancyfoot[C]{\thepage}
\fancypagestyle{plain}{\fancyhf{}\fancyhead[L]{\small HPS Gaussian-Process Resonance Search}
\fancyhead[R]{\small v6.2 / Signal recovery}\fancyfoot[C]{\thepage}}
\setlength{\emergencystretch}{2em}\setlength{\tabcolsep}{4pt}
\renewcommand{\arraystretch}{1.08}
\newcommand{\fig}[3]{\begin{figure}[H]\centering
\includegraphics[width=.98\linewidth,height=#1\textheight,keepaspectratio]{#2}\caption{#3}\end{figure}}
% Generated from saved CSV ledgers; do not edit numerical values here.
\newcommand{\ToyCount}{40}
\newcommand{\TotalFitRows}{12,320}
\newcommand{\PrimaryFitRows}{3,520}
\newcommand{\UniqueInjectedToys}{1,760}
\newcommand{\UniqueNullToys}{440}
\newcommand{\IncPoleRange}{0.715--0.895}
\newcommand{\IncCoreRange}{0.780--0.938}
\newcommand{\RawPoleRange}{0.822--1.011}
\newcommand{\RawCoreRange}{0.845--1.073}
\newcommand{\CleanRange}{0.966--1.158}
\newcommand{\KnownRange}{0.928--1.022}
\newcommand{\PullMeanRange}{-2.26--0.48}
\newcommand{\PullWidthRange}{0.79--1.12}
\newcommand{\ContainmentMin}{15}
\newcommand{\ContainmentMax}{40}
\newcommand{\ContainmentStep}{2.5}
\newcommand{\WidthRelativeSE}{11}
\newcommand{\ExamplePoleYield}{24,649}
\newcommand{\ExamplePoleRecovery}{0.822}
\newcommand{\ExamplePolePull}{-2.04}
\newcommand{\ExamplePoleInclusion}{15}
\newcommand{\ExamplePoleInclusionInterval}{0.23--0.54}
\newcommand{\ExampleCoreYield}{25,360}
\newcommand{\ExampleCoreRecovery}{0.845}
\newcommand{\ExampleCorePull}{-1.72}
\newcommand{\ExampleCoreInclusion}{22}
\newcommand{\ExampleCoreInclusionInterval}{0.38--0.71}
\newcommand{\AsimovInjected}{24,023/24,859}
\newcommand{\AsimovClean}{29,743/29,724}
\newcommand{\AsimovKnown}{30,000/30,000}
\newcommand{\IncrementStabilityPercent}{0.29}

\title{HPS Gaussian-Process Resonance Search\\[.5em]
\large Reconstructed-MC Injection and Signal Recovery\\Standalone Study v6.2}
\author{Emrys Peets\\Stanford University and SLAC National Accelerator Laboratory}
\date{23 September 2026}
\begin{document}
\maketitle
\begin{abstract}
Signal extraction is tested by adding exactly 1,000, 5,000, 10,000 and
30,000 selected candidates from native smeared 2021 MC distributions to
\ToyCount{} background toys at each mass from 60 through 260~MeV in
20~MeV steps. Paired fits compare windows centered at the generated mass
and the reconstructed MC core, with the same nominal width. At 30,000
injected candidates, subtracting the paired zero-injection fit gives a
mean recovered increment of \IncPoleRange{} of the injected yield for
the original window and \IncCoreRange{} for the shifted window over
60--240~MeV. The shift reduces the loss but does not restore full
recovery. Controls with signal removed from the training sidebands
attribute much of the deficit to the GP response to the injected tails.
Pulls and profile likelihood inclusion counts accompany the yield tests.
\end{abstract}

\section{Result and relation to v6.1}
The selected-candidate MC distributions examined in v6.1 have broad
tails. Injected candidates outside the fitted window enter the sideband
sample and can change the GP prediction beneath the signal core. The
recovery tests include this response to the injected counts.

\begin{table}[H]\centering\small
\begin{tabular}{rrr}
\toprule
Injected $N$ & Pole-centered & MC-core-centered \\
\midrule
1,000 & -0.089--5.227 & -0.069--5.566 \\
5,000 & 0.689--1.737 & 0.740--1.820 \\
10,000 & 0.783--1.301 & 0.837--1.372 \\
30,000 & 0.822--1.011 & 0.845--1.073 \\
\bottomrule
\end{tabular}

\caption{Ranges of the mean fitted yield divided by the injected yield
over 60--240~MeV. These summarize the ten mass points; they are not
confidence intervals. The low-injection ratios include the background
offset and should be read with their toy errors in Figure~\ref{fig:recovery}.}
\end{table}

The tests use fixed empirical MC shapes and one background mean. The
v6.1 audit did not establish signal-daughter association or full
equivalence between the MC and data selections. The template and its
MC-only center are held fixed; finite-MC, background-truth and
detector-response uncertainties are not sampled.

\clearpage
\section{Toy ensemble and yield normalization}
The GP arithmetic mean fitted over the full mass support of the 2021
10\% spectrum in v5.8.2 defines the background expectation $\mu_i$ over the unchanged analysis
support $D=[36,299.75]$~MeV. Independent Poisson spectra are generated
for each mass and toy index. At a fixed mass, the same spectrum is used
at all four injection levels and for the zero-injection fit. The two
window prescriptions are fitted to the same injected spectrum.

The signal probabilities $p_i$ are obtained from the native selected-MC
histogram at that generated mass. They use the full selected normalization,
including the recorded overflow. Counts are drawn as
\begin{equation}
 B_i\sim\operatorname{Poisson}(\mu_i),\qquad
 (S_1,\ldots,S_K,S_{\mathrm{below}},S_{\mathrm{above}})
 \sim\operatorname{Multinomial}(N;\bm p),\qquad n_i=B_i+S_i.
 \label{eq:toys}
\end{equation}
The last two categories retain candidates outside $D$, so the total
draw is exactly $N$ even when fewer candidates enter the analysis.
CDF rebinning treats the density as uniform within each source histogram
bin. No additional smearing, translation, truth selection or interpolation
between generated masses is applied. Here and below, ``yield'' means
selected-candidate yield, not a number of generated resonances.

The signal contribution in a fitted bin is $A p_i$. In particular,
\begin{equation}
 f_D=\sum_{i\in D}p_i,\qquad f_W=\sum_{i\in W}p_i
 \label{eq:fractions}
\end{equation}
remain the probabilities for a selected candidate to enter $D$ and $W$.
The template is normalized before either restriction, so $A$ estimates
the full selected yield $N$. Expected fractions and
actual draws inside the support, window and training bins are saved for
each toy.

\begin{table}[H]\centering\small
\begin{tabular}{rrrrrr}
\toprule
$m_0$ & $c(m_0)$ & $\sigma_m$ & $f_D$ & $f_{W,\mathrm{P}}$ & $f_{W,\mathrm{C}}$ \\
\midrule
60 & 58.792 & 2.075 & 0.998 & 0.392 & 0.396 \\
80 & 78.103 & 2.288 & 0.999 & 0.675 & 0.709 \\
100 & 97.688 & 2.570 & 0.999 & 0.731 & 0.781 \\
120 & 117.319 & 2.920 & 0.999 & 0.745 & 0.803 \\
140 & 136.956 & 3.339 & 0.999 & 0.764 & 0.812 \\
160 & 156.658 & 3.827 & 0.999 & 0.751 & 0.816 \\
180 & 176.393 & 4.383 & 0.999 & 0.752 & 0.803 \\
200 & 196.182 & 5.008 & 0.998 & 0.756 & 0.792 \\
220 & 215.941 & 5.702 & 0.998 & 0.737 & 0.771 \\
240 & 235.617 & 6.465 & 0.997 & 0.722 & 0.748 \\
260$^{*}$ & 255.173 & 7.296 & 0.995 & 0.701 & 0.722 \\
\bottomrule
\end{tabular}

\caption{Fixed MC centers, nominal resolutions and selected fractions.
Masses, centers and resolutions are in MeV. P and C denote the pole- and
MC-core-centered masks. The asterisk marks the 260~MeV extension, which
uses the archived 250~MeV kernel state.}
\end{table}

\clearpage
\section{Extraction and uncertainty calculation}
The pole-centered mask is $m_0\pm2.25\sigma_m(m_0)$; the shifted mask is
$c(m_0)\pm2.25\sigma_m(m_0)$. The center $c$ is determined from MC alone
using the local integrated Gaussian plus nonnegative affine pedestal
fit defined in v6.1. The same definition is applied at 260~MeV.
The MC shape does not move when the window moves. The resolution is
\begin{equation}
 \sigma_m(m)=0.00184825-0.001375m+0.085875m^2,
 \qquad m,\sigma_m\ \hbox{in GeV}.
\end{equation}
Kernel parameters remain anchored at $m_0$, with the archived 250~MeV
state used for the 260~MeV extension. The MC center and nominal width
are fixed before drawing the toys.

In the primary fit the GP is trained on the injected toy's exterior
bins. Its archived kernel hyperparameters remain fixed, while the
count-dependent targets, noise terms, conditional mean $\widehat{\bm b}$
and predictive covariance are recomputed. The v6.1 factorization retains
the covariance modes represented by $L$, with $V=LL^T$. The inherited
likelihood profiles their nuisance coordinates $\bm\theta$:
\begin{equation}
 \begin{aligned}
 -\ln\mathcal{L}(A,\bm\theta)&=
 \sum_{i\in W}\left[\lambda_i-n_i\ln\lambda_i\right]
 +\frac12\bm\theta^T\bm\theta,\\
 \lambda_i&=\widehat b_i+(L\bm\theta)_i+A p_i>0.
 \end{aligned}
 \label{eq:likelihood}
\end{equation}
The fitted amplitude $\widehat A$ is allowed to be negative, retaining
the signed behavior around zero signal. Its uncertainty $\sigma_A$ is
obtained from the observed curvature of the profiled likelihood.

Two controls use the same injected counts in the fitted window.
The clean-sideband control trains the GP on the background-only bins
of that toy, removing the injected tail from training. The
known-background control fixes the background to $\mu_i$ and has no GP
nuisance parameters.
The former isolates the change caused by signal in the training sample;
the latter also removes background-estimation fluctuations and the GP
nuisance covariance.

Writing $\ell_{\mathrm{prof}}(A)$ for the negative log likelihood minimized
over $\bm\theta$ at fixed $A$, the pull and likelihood ratio at the
injected truth are
\begin{equation}
 p=\frac{\widehat A-N}{\sigma_A},\qquad
 q_N=2\left[\ell_{\mathrm{prof}}(N)-
                  \ell_{\mathrm{prof}}(\widehat A)\right].
 \label{eq:pull}
\end{equation}
The nominal 68.27\% and 95\% profile sets include $N$ when $q_N\leq1$
and $q_N\leq\chi^2_1(0.95)$, respectively. The fraction of toys satisfying
each condition is reported with an exact 95\% Clopper--Pearson interval.
Because the signal in
Eq.~\eqref{eq:toys} has an exact multinomial total while
Eq.~\eqref{eq:likelihood} uses Poisson counts, width one and the nominal
inclusion fractions are reference values, including for known background.

\clearpage
\section{Mean recovered yield}
For each mass, level and method, the raw recovery ratio over
$T=\ToyCount$ toys is
\begin{equation}
 R(m_0,N)=\frac{1}{NT}
              \sum_{t=1}^{T}\widehat A_t(m_0,N).
\end{equation}
It includes the response to the injected signal and any fitted offset
already present at zero injection. At $N=30{,}000$, the ranges over
60--240~MeV are \RawPoleRange{} for the pole-centered window and
\RawCoreRange{} for the shifted window.

\fig{.60}{../figures/recovery.pdf}{Mean raw recovery from the primary
injected-sideband fits. Error bars are the sample standard error across
\ToyCount{} toys; the horizontal line is exact mean recovery.
The gray band marks the separate 260~MeV extension.\label{fig:recovery}}

At the smaller injections, the fitted background offset and its
fluctuations can be comparable to $N$, producing large changes in the
raw ratio. Section~\ref{sec:increment} subtracts each toy's zero-injection
fit to isolate the incremental response.

\clearpage
\section{Pulls and profile inclusion}
The mean pull expresses the average residual in fitted uncertainty
units. A zero mean pull need not imply $E[\widehat A]=N$, since the
uncertainty $\sigma_A$ can vary with the residual. The tables therefore
give the mean yield and mean uncertainty alongside the pull summaries.

\fig{.61}{../figures/pull_mean.pdf}{Mean pulls in the primary fits.
Bars are the sample pull standard deviation divided by
$\sqrt{\ToyCount}$. Both methods share the same background and injection
draws.}

At 240~MeV and $N=30{,}000$, the mean fitted yields are
\ExamplePoleYield{} and \ExampleCoreYield{} candidates for the pole- and
core-centered windows. The corresponding mean pulls are
\ExamplePolePull{} and \ExampleCorePull{}. The shift reduces the
negative residual but leaves a substantial deficit.

\clearpage
\subsection{Dispersion and interval inclusion}
The pull width is the sample standard deviation, calculated with the
$T-1$ denominator. It is estimated directly from the individual
pulls without fitting a Gaussian to the histograms. With \ToyCount{}
toys its sampling uncertainty remains appreciable: even for Gaussian
pulls, the width has an approximate \WidthRelativeSE\% relative standard
error. The four injection levels also share background spectra, so their
summaries are correlated.

\fig{.57}{../figures/pull_width.pdf}{Sample pull widths for the primary
fits. The dashed line marks one. The exact-$N$ signal ensemble and the
GP uncertainty treatment need not produce this width even in the
absence of mean bias.}

At 240~MeV and $N=30{,}000$, the nominal 95\% profile set contains the
truth in \ExamplePoleInclusion{}/\ToyCount{} pole-centered fits and
\ExampleCoreInclusion{}/\ToyCount{} shifted fits. The corresponding
95\% binomial intervals are \ExamplePoleInclusionInterval{} and
\ExampleCoreInclusionInterval{}. These counts describe this toy
ensemble; \ToyCount{} trials do not establish calibrated coverage.

\clearpage
\section{Zero-injection behavior}
The $N=0$ fits use the same background spectra that appear at the four
nonzero injection levels. They reveal the baseline fluctuations and
offsets that contribute to $R$ at small $N$. Clean and injected GP
sidebands coincide at zero injection, so only one GP result is needed.

\fig{.62}{../figures/null_bias.pdf}{Mean signed fitted yield at zero
injection, with standard errors across \ToyCount{} background toys.
The upper panel uses GP background estimation; the lower panel fixes
the known generating expectation.}

The uncertainty in the known-background control assumes that the
background mean is known exactly. It excludes the statistical cost of
estimating that mean from sideband data and cannot be substituted for
the uncertainty in the GP extraction. The paired subtraction below uses
each method's own zero-injection fit.

\clearpage
\section{Signal in the training sidebands}
The 30,000-candidate injections make the effect of the broad MC tails
visible relative to the statistical fluctuations. Across both masks
and 60--240~MeV, clean-sideband mean recovery spans \CleanRange{}, while
known-background recovery spans \KnownRange{}.

\fig{.57}{../figures/controls_30000.pdf}{Recovery and mean pulls at
$N=30{,}000$ for the three extraction conditions. Only the training
sample changes between the injected- and clean-sideband GP fits. The
known-background fit additionally fixes the fitted background vector.}

The effect also occurs without toy fluctuations. For a deterministic
mean-count spectrum at 240~MeV with $N=30{,}000$, the extracted yields
are \AsimovInjected{} with injected sidebands,
\AsimovClean{} with clean sidebands and \AsimovKnown{} with known
background, where each pair lists pole-centered/core-centered results.
The additional sideband counts raise the background prediction in the
signal window, leaving less yield for the signal component. Centering
the mask on the MC core reduces this loss but does not remove the broad
tail from the training sample.

\clearpage
\section{Paired response to the injected signal}\label{sec:increment}
To separate the response to additional signal from the zero-injection
offset, define
\begin{equation}
 R_{\mathrm{inc}}(m_0,N)=\frac{1}{NT}
 \sum_{t=1}^{T}
 \left[\widehat A_t(m_0,N)-\widehat A_t(m_0,0)\right].
 \label{eq:increment}
\end{equation}
This quantity measures the change in fitted yield caused by adding
signal to the same background toy. The raw estimator and its pulls
retain the zero-injection offset.

\fig{.53}{../figures/paired_comparison.pdf}{The upper panel compares
the two fitted yields within each identical toy, divided by $N$.
Error bars use the standard error of that paired difference. Lower
panels show $R_{\mathrm{inc}}$ for each window prescription.}

At $N=30{,}000$, $R_{\mathrm{inc}}$ spans \IncPoleRange{} for the
pole-centered mask and \IncCoreRange{} for the shifted mask over
60--240~MeV. Adding the second set of 20 toys changes this increment by
at most \IncrementStabilityPercent{} percentage points across those
masses and methods. The first 20 draws and their fit results are
unchanged and are preserved in the archived 20-toy release. The ensemble
contains \UniqueInjectedToys{} injected toys and \UniqueNullToys{}
background toys, giving \PrimaryFitRows{} primary nonzero-injection
fits and \TotalFitRows{} fits including the controls. The persistent
deficit in the paired response agrees with the clean-sideband and
mean-count comparisons.

\clearpage
\appendix
\section{Per-mass pull distributions}
Each outline below contains \ToyCount{} primary-fit pulls. The methods
share bin edges within a mass panel, while ranges may differ between
masses. Histograms show counts rather than a smooth density estimate.

\fig{.78}{../figures/pull_gallery_N01000.pdf}{Pulls for $N=1{,}000$.
The shaded 260~MeV panel denotes the extraction extension.}

\clearpage
\subsection{Five thousand injected candidates}
\fig{.83}{../figures/pull_gallery_N05000.pdf}{Pulls for $N=5{,}000$;
\ToyCount{} toys per method and mass.}

\clearpage
\subsection{Ten thousand injected candidates}
\fig{.83}{../figures/pull_gallery_N10000.pdf}{Pulls for $N=10{,}000$;
\ToyCount{} toys per method and mass.}

\clearpage
\subsection{Thirty thousand injected candidates}
\fig{.83}{../figures/pull_gallery_N30000.pdf}{Pulls for $N=30{,}000$;
\ToyCount{} toys per method and mass.}

\clearpage
\section{Primary-fit numerical results}
Tables below contain every primary mass/injection/method cell. P denotes
the pole-centered window and C the MC-core-centered window. The mass
$m_0$ is in MeV; $\overline{\widehat A}$ and
$\overline{\sigma_A}$ are in selected candidates. The raw recovery is
$R$, and $\overline p$ and $s_p$ are the mean and sample standard
deviation of the pulls. The last two columns give the nominal 68.27\%
and 95\% profile-inclusion counts out of \ToyCount{}.

\begin{table}[H]\centering\small
\begin{tabular}{rlrrrrrrr}
\toprule
$m_0$ & Mask & $\overline{\widehat A}$ & $\overline{\sigma_A}$ & $R$ & $\overline p$ & $s_p$ & $k_{68}$ & $k_{95}$ \\
\midrule
60 & P & 4,900 & 19,376 & 4.900 & 0.20 & 1.03 & 24 & 38 \\
60 & C & 5,566 & 21,052 & 5.566 & 0.22 & 0.98 & 27 & 38 \\
80 & P & 5,227 & 8,861 & 5.227 & 0.48 & 0.95 & 26 & 36 \\
80 & C & 5,411 & 9,374 & 5.411 & 0.47 & 0.92 & 28 & 36 \\
100 & P & 300 & 6,983 & 0.300 & -0.10 & 0.89 & 33 & 38 \\
100 & C & 641 & 7,360 & 0.641 & -0.05 & 0.90 & 29 & 38 \\
120 & P & 1,475 & 5,708 & 1.475 & 0.08 & 1.02 & 29 & 37 \\
120 & C & 1,667 & 5,985 & 1.667 & 0.11 & 0.96 & 29 & 38 \\
140 & P & 1,158 & 4,761 & 1.158 & 0.03 & 1.11 & 29 & 37 \\
140 & C & 1,361 & 4,942 & 1.361 & 0.07 & 1.07 & 29 & 37 \\
160 & P & -89 & 3,950 & -0.089 & -0.28 & 0.96 & 27 & 38 \\
160 & C & -69 & 4,147 & -0.069 & -0.26 & 0.92 & 29 & 38 \\
180 & P & 698 & 3,412 & 0.698 & -0.09 & 0.93 & 29 & 38 \\
180 & C & 861 & 3,549 & 0.861 & -0.04 & 0.86 & 30 & 40 \\
200 & P & 335 & 3,041 & 0.335 & -0.22 & 0.86 & 31 & 38 \\
200 & C & 490 & 3,132 & 0.490 & -0.16 & 0.79 & 32 & 39 \\
220 & P & 1,335 & 2,752 & 1.335 & 0.12 & 1.00 & 29 & 38 \\
220 & C & 1,317 & 2,833 & 1.317 & 0.11 & 1.01 & 27 & 38 \\
240 & P & 1,202 & 2,596 & 1.202 & 0.08 & 0.96 & 30 & 38 \\
240 & C & 1,097 & 2,669 & 1.097 & 0.04 & 0.93 & 29 & 38 \\
260$^{*}$ & P & 1,752 & 2,659 & 1.752 & 0.28 & 0.98 & 25 & 38 \\
260$^{*}$ & C & 1,777 & 2,745 & 1.777 & 0.28 & 0.88 & 27 & 37 \\
\bottomrule
\end{tabular}

\caption{Primary fits for $N=1{,}000$. The asterisk marks 260~MeV.}
\end{table}

\noindent The 260~MeV extension uses the archived 250~MeV kernel and
is excluded from the 60--240~MeV summary ranges. Full-precision results,
all controls, Monte Carlo standard errors and Clopper--Pearson intervals
are retained in \path{results/summary.csv}; individual fits are in
\path{results/toys.csv}. Input identities and the fixed construction
are recorded in \path{protocol.json} and the provenance files.

\clearpage
\subsection{Five thousand injected candidates}
\begin{table}[H]\centering\small
\begin{tabular}{rlrrrrrrr}
\toprule
$m_0$ & Mask & $\overline{\widehat A}$ & $\overline{\sigma_A}$ & $R$ & $\overline p$ & $s_p$ & $k_{68}$ & $k_{95}$ \\
\midrule
60 & P & 7,749 & 19,377 & 1.550 & 0.14 & 1.02 & 25 & 37 \\
60 & C & 8,678 & 21,054 & 1.736 & 0.18 & 0.98 & 28 & 38 \\
80 & P & 8,686 & 8,861 & 1.737 & 0.42 & 0.95 & 26 & 36 \\
80 & C & 9,101 & 9,375 & 1.820 & 0.44 & 0.92 & 28 & 36 \\
100 & P & 3,822 & 6,984 & 0.764 & -0.17 & 0.89 & 33 & 38 \\
100 & C & 4,394 & 7,361 & 0.879 & -0.08 & 0.90 & 29 & 38 \\
120 & P & 5,019 & 5,709 & 1.004 & 0.00 & 1.02 & 31 & 37 \\
120 & C & 5,437 & 5,986 & 1.087 & 0.07 & 0.96 & 30 & 38 \\
140 & P & 4,726 & 4,762 & 0.945 & -0.06 & 1.11 & 30 & 37 \\
140 & C & 5,105 & 4,943 & 1.021 & 0.02 & 1.07 & 30 & 37 \\
160 & P & 3,446 & 3,951 & 0.689 & -0.39 & 0.96 & 26 & 38 \\
160 & C & 3,698 & 4,148 & 0.740 & -0.31 & 0.91 & 30 & 38 \\
180 & P & 4,183 & 3,413 & 0.837 & -0.24 & 0.93 & 29 & 39 \\
180 & C & 4,531 & 3,550 & 0.906 & -0.13 & 0.86 & 29 & 39 \\
200 & P & 3,845 & 3,043 & 0.769 & -0.38 & 0.86 & 28 & 38 \\
200 & C & 4,139 & 3,134 & 0.828 & -0.27 & 0.79 & 31 & 39 \\
220 & P & 4,707 & 2,755 & 0.941 & -0.11 & 1.00 & 28 & 37 \\
220 & C & 4,817 & 2,835 & 0.963 & -0.06 & 1.01 & 26 & 39 \\
240 & P & 4,455 & 2,598 & 0.891 & -0.21 & 0.95 & 29 & 38 \\
240 & C & 4,459 & 2,672 & 0.892 & -0.20 & 0.93 & 30 & 38 \\
260$^{*}$ & P & 4,839 & 2,664 & 0.968 & -0.06 & 0.98 & 28 & 38 \\
260$^{*}$ & C & 4,977 & 2,750 & 0.995 & -0.01 & 0.88 & 28 & 40 \\
\bottomrule
\end{tabular}

\caption{Primary fits for $N=5{,}000$; \ToyCount{} toys per row.
Column definitions are given at the start of Appendix~B.}
\end{table}

\clearpage
\subsection{Ten thousand injected candidates}
\begin{table}[H]\centering\small
\begin{tabular}{rlrrrrrrr}
\toprule
$m_0$ & Mask & $\overline{\widehat A}$ & $\overline{\sigma_A}$ & $R$ & $\overline p$ & $s_p$ & $k_{68}$ & $k_{95}$ \\
\midrule
60 & P & 11,448 & 19,379 & 1.145 & 0.08 & 1.02 & 25 & 38 \\
60 & C & 12,682 & 21,056 & 1.268 & 0.13 & 0.98 & 27 & 38 \\
80 & P & 13,014 & 8,862 & 1.301 & 0.34 & 0.95 & 27 & 36 \\
80 & C & 13,724 & 9,376 & 1.372 & 0.40 & 0.92 & 28 & 38 \\
100 & P & 8,261 & 6,985 & 0.826 & -0.25 & 0.89 & 30 & 39 \\
100 & C & 9,125 & 7,362 & 0.913 & -0.12 & 0.90 & 30 & 38 \\
120 & P & 9,429 & 5,710 & 0.943 & -0.10 & 1.02 & 30 & 37 \\
120 & C & 10,133 & 5,987 & 1.013 & 0.02 & 0.96 & 31 & 38 \\
140 & P & 9,206 & 4,763 & 0.921 & -0.17 & 1.12 & 29 & 36 \\
140 & C & 9,805 & 4,944 & 0.981 & -0.04 & 1.08 & 30 & 37 \\
160 & P & 7,830 & 3,952 & 0.783 & -0.55 & 0.96 & 22 & 39 \\
160 & C & 8,367 & 4,149 & 0.837 & -0.39 & 0.91 & 28 & 38 \\
180 & P & 8,548 & 3,415 & 0.855 & -0.43 & 0.94 & 27 & 39 \\
180 & C & 9,137 & 3,552 & 0.914 & -0.24 & 0.87 & 28 & 39 \\
200 & P & 8,199 & 3,045 & 0.820 & -0.59 & 0.86 & 28 & 38 \\
200 & C & 8,659 & 3,136 & 0.866 & -0.43 & 0.80 & 31 & 39 \\
220 & P & 8,912 & 2,757 & 0.891 & -0.39 & 0.99 & 27 & 37 \\
220 & C & 9,196 & 2,838 & 0.920 & -0.28 & 0.99 & 28 & 38 \\
240 & P & 8,454 & 2,602 & 0.845 & -0.59 & 0.97 & 28 & 38 \\
240 & C & 8,616 & 2,676 & 0.862 & -0.52 & 0.94 & 27 & 38 \\
260$^{*}$ & P & 8,629 & 2,669 & 0.863 & -0.51 & 0.99 & 24 & 36 \\
260$^{*}$ & C & 8,896 & 2,755 & 0.890 & -0.40 & 0.90 & 29 & 38 \\
\bottomrule
\end{tabular}

\caption{Primary fits for $N=10{,}000$; \ToyCount{} toys per row.
Column definitions are given at the start of Appendix~B.}
\end{table}

\clearpage
\subsection{Thirty thousand injected candidates}
\begin{table}[H]\centering\small
\begin{tabular}{rlrrrrrrr}
\toprule
$m_0$ & Mask & $\overline{\widehat A}$ & $\overline{\sigma_A}$ & $R$ & $\overline p$ & $s_p$ & $k_{68}$ & $k_{95}$ \\
\midrule
60 & P & 25,657 & 19,387 & 0.855 & -0.22 & 1.02 & 28 & 38 \\
60 & C & 28,206 & 21,064 & 0.940 & -0.08 & 0.97 & 28 & 38 \\
80 & P & 30,320 & 8,866 & 1.011 & 0.04 & 0.95 & 27 & 40 \\
80 & C & 32,185 & 9,379 & 1.073 & 0.23 & 0.92 & 29 & 39 \\
100 & P & 25,830 & 6,989 & 0.861 & -0.60 & 0.89 & 27 & 39 \\
100 & C & 27,846 & 7,365 & 0.928 & -0.29 & 0.90 & 30 & 39 \\
120 & P & 27,046 & 5,713 & 0.902 & -0.52 & 1.02 & 27 & 37 \\
120 & C & 28,857 & 5,991 & 0.962 & -0.19 & 0.96 & 29 & 38 \\
140 & P & 27,100 & 4,768 & 0.903 & -0.61 & 1.12 & 21 & 36 \\
140 & C & 28,571 & 4,949 & 0.952 & -0.29 & 1.07 & 26 & 36 \\
160 & P & 25,397 & 3,958 & 0.847 & -1.16 & 0.94 & 13 & 33 \\
160 & C & 27,097 & 4,155 & 0.903 & -0.70 & 0.90 & 21 & 38 \\
180 & P & 26,085 & 3,422 & 0.870 & -1.14 & 0.94 & 19 & 32 \\
180 & C & 27,617 & 3,559 & 0.921 & -0.67 & 0.87 & 25 & 38 \\
200 & P & 25,582 & 3,054 & 0.853 & -1.45 & 0.86 & 11 & 29 \\
200 & C & 26,738 & 3,144 & 0.891 & -1.04 & 0.79 & 18 & 37 \\
220 & P & 25,741 & 2,768 & 0.858 & -1.54 & 1.01 & 9 & 28 \\
220 & C & 26,677 & 2,849 & 0.889 & -1.17 & 1.02 & 20 & 32 \\
240 & P & 24,649 & 2,616 & 0.822 & -2.04 & 0.94 & 6 & 15 \\
240 & C & 25,360 & 2,690 & 0.845 & -1.72 & 0.91 & 8 & 22 \\
260$^{*}$ & P & 23,918 & 2,692 & 0.797 & -2.26 & 0.95 & 4 & 15 \\
260$^{*}$ & C & 24,795 & 2,778 & 0.826 & -1.87 & 0.87 & 5 & 21 \\
\bottomrule
\end{tabular}

\caption{Primary fits for $N=30{,}000$; \ToyCount{} toys per row.
Column definitions are given at the start of Appendix~B.}
\end{table}
\end{document}
